
\documentclass{article}
\usepackage{colortbl}
\usepackage{makecell}
\usepackage{multirow}
\usepackage{supertabular}

\begin{document}

\newcounter{utterance}

\centering \large Interaction Transcript for game `matchit\_ascii', experiment `same\_grid', episode 3 with glm{-}5.2.
\vspace{24pt}

{ \footnotesize  \setcounter{utterance}{1}
\setlength{\tabcolsep}{0pt}
\begin{supertabular}{c@{$\;$}|p{.15\linewidth}@{}p{.15\linewidth}p{.15\linewidth}p{.15\linewidth}p{.15\linewidth}p{.15\linewidth}}
    \# & $\;$A & \multicolumn{4}{c}{Game Master} & $\;\:$B\\
    \hline

    \theutterance \stepcounter{utterance}  
    & & \multicolumn{4}{p{0.6\linewidth}}{
        \cellcolor[rgb]{0.9,0.9,0.9}{
            \makecell[{{p{\linewidth}}}]{
                \texttt{\tiny{[P1$\langle$GM]}}
                \texttt{Vous participez à un jeu de devinettes collaboratif.} \\
\texttt{L'objectif est de déterminer si cette grille et une autre image que je suis le seul à voir sont identiques.} \\
\texttt{Veuillez d'abord décrire votre grille. Ensuite, je fournirai ma description et nous pourrons nous poser mutuellement des questions sur les grilles pour déterminer si elles sont identiques. Commencez maintenant votre brève description de la grille par "DESCRIPTION :" suivi de la description. N'ajoutez rien d'autre. Voici votre grille :} \\
\texttt{X X X X X} \\
\texttt{▢ ▢ ▢ ▢ X} \\
\texttt{▢ ▢ X X X} \\
\texttt{▢ ▢ ▢ ▢ X} \\
\texttt{▢ ▢ ▢ ▢ X} \\
            }
        }
    }
    & \\ \\

    \theutterance \stepcounter{utterance}  
    & \multicolumn{4}{p{0.6\linewidth}}{
        \cellcolor[rgb]{0.9,0.9,0.9}{
            \makecell[{{p{\linewidth}}}]{
                \texttt{\tiny{[P1$\rangle$GM]}}
                \texttt{DESCRIPTION : Grille de 5 lignes par 5 colonnes. Ligne 1 : cinq X. Ligne 2 : quatre ▢ puis un X. Ligne 3 : deux ▢ puis trois X. Ligne 4 : quatre ▢ puis un X. Ligne 5 : quatre ▢ puis un X.} \\
            }
        }
    }
    & & \\ \\

    \theutterance \stepcounter{utterance}  
    & & & \multicolumn{2}{p{0.3\linewidth}}{
        \cellcolor[rgb]{0.9,0.9,0.9}{
            \makecell[{{p{\linewidth}}}]{
                \texttt{\tiny{[GM$|$GM]}}
                \texttt{continue} \\
            }
        }
    }
    & & \\ \\

    \theutterance \stepcounter{utterance}  
    & & \multicolumn{4}{p{0.6\linewidth}}{
        \cellcolor[rgb]{0.9,0.9,0.9}{
            \makecell[{{p{\linewidth}}}]{
                \texttt{\tiny{[GM$\rangle$P2]}}
                \texttt{Vous participez à un jeu de devinettes collaboratif.} \\
\texttt{L'objectif est de déterminer si cette grille et une autre image que je suis le seul à voir sont identiques.} \\
\texttt{Veuillez d'abord décrire votre grille. Ensuite, je fournirai ma description et nous pourrons nous poser mutuellement des questions sur les grilles pour déterminer si elles sont identiques. Commencez maintenant votre brève description de la grille par "DESCRIPTION :" suivi de la description. N'ajoutez rien d'autre. Voici votre grille :} \\
\texttt{X X X X X} \\
\texttt{▢ ▢ ▢ ▢ X} \\
\texttt{▢ ▢ X X X} \\
\texttt{▢ ▢ ▢ ▢ X} \\
\texttt{▢ ▢ ▢ ▢ X} \\
            }
        }
    }
    & \\ \\

    \theutterance \stepcounter{utterance}  
    & & & \multicolumn{4}{p{0.6\linewidth}}{
        \cellcolor[rgb]{0.9,0.9,0.9}{
            \makecell[{{p{\linewidth}}}]{
                \texttt{\tiny{[GM$\langle$P2]}}
                \texttt{DESCRIPTION : La grille est composée de 5 lignes et 5 colonnes. La première ligne est entièrement remplie de "X". Les quatre autres lignes commencent par deux cases vides "▢" suivies de deux "X" pour la troisième ligne et d'un seul "X" pour les deuxième, quatrième et cinquième lignes.} \\
            }
        }
    }
     \\ \\

    \theutterance \stepcounter{utterance}  
    & & & \multicolumn{2}{p{0.3\linewidth}}{
        \cellcolor[rgb]{0.9,0.9,0.9}{
            \makecell[{{p{\linewidth}}}]{
                \texttt{\tiny{[GM$|$GM]}}
                \texttt{continue} \\
            }
        }
    }
    & & \\ \\

    \theutterance \stepcounter{utterance}  
    & & \multicolumn{4}{p{0.6\linewidth}}{
        \cellcolor[rgb]{0.9,0.9,0.9}{
            \makecell[{{p{\linewidth}}}]{
                \texttt{\tiny{[GM$\rangle$P2]}}
                \texttt{Voici ma DESCRIPTIONDESCRIPTION : Grille de 5 lignes par 5 colonnes. Ligne 1 : cinq X. Ligne 2 : quatre ▢ puis un X. Ligne 3 : deux ▢ puis trois X. Ligne 4 : quatre ▢ puis un X. Ligne 5 : quatre ▢ puis un X.} \\
\texttt{Posez maintenant une question afin de découvrir de nouveaux aspects de ma grille qui pourraient être différents de votre grille. Commencez par "QUESTION :" et n'ajoutez rien d'autre.} \\
            }
        }
    }
    & \\ \\

    \theutterance \stepcounter{utterance}  
    & & & \multicolumn{4}{p{0.6\linewidth}}{
        \cellcolor[rgb]{0.9,0.9,0.9}{
            \makecell[{{p{\linewidth}}}]{
                \texttt{\tiny{[GM$\langle$P2]}}
                \texttt{QUESTION : Dans la deuxième ligne de votre grille, les quatre cases vides "▢" sont{-}elles suivies du "X" en cinquième position (en bordure droite) ?} \\
            }
        }
    }
     \\ \\

    \theutterance \stepcounter{utterance}  
    & & & \multicolumn{2}{p{0.3\linewidth}}{
        \cellcolor[rgb]{0.9,0.9,0.9}{
            \makecell[{{p{\linewidth}}}]{
                \texttt{\tiny{[GM$|$GM]}}
                \texttt{continue} \\
            }
        }
    }
    & & \\ \\

    \theutterance \stepcounter{utterance}  
    & & \multicolumn{4}{p{0.6\linewidth}}{
        \cellcolor[rgb]{0.9,0.9,0.9}{
            \makecell[{{p{\linewidth}}}]{
                \texttt{\tiny{[P1$\langle$GM]}}
                \texttt{Voici ma DESCRIPTIONDESCRIPTION : La grille est composée de 5 lignes et 5 colonnes. La première ligne est entièrement remplie de "X". Les quatre autres lignes commencent par deux cases vides "▢" suivies de deux "X" pour la troisième ligne et d'un seul "X" pour les deuxième, quatrième et cinquième lignes.} \\
\texttt{QUESTION : Dans la deuxième ligne de votre grille, les quatre cases vides "▢" sont{-}elles suivies du "X" en cinquième position (en bordure droite) ?Commencez votre réponse par "RÉPONSE :" et n'ajoutez rien d'autre.} \\
            }
        }
    }
    & \\ \\

    \theutterance \stepcounter{utterance}  
    & \multicolumn{4}{p{0.6\linewidth}}{
        \cellcolor[rgb]{0.9,0.9,0.9}{
            \makecell[{{p{\linewidth}}}]{
                \texttt{\tiny{[P1$\rangle$GM]}}
                \texttt{RÉPONSE : Oui, dans la deuxième ligne, les quatre cases vides "▢" sont suivies d'un "X" en cinquième position (bordure droite).} \\
            }
        }
    }
    & & \\ \\

    \theutterance \stepcounter{utterance}  
    & & & \multicolumn{2}{p{0.3\linewidth}}{
        \cellcolor[rgb]{0.9,0.9,0.9}{
            \makecell[{{p{\linewidth}}}]{
                \texttt{\tiny{[GM$|$GM]}}
                \texttt{continue} \\
            }
        }
    }
    & & \\ \\

    \theutterance \stepcounter{utterance}  
    & & \multicolumn{4}{p{0.6\linewidth}}{
        \cellcolor[rgb]{0.9,0.9,0.9}{
            \makecell[{{p{\linewidth}}}]{
                \texttt{\tiny{[P1$\langle$GM]}}
                \texttt{Posez maintenant une question afin de découvrir de nouveaux aspects de ma grille qui pourraient être différents de votre grille. Commencez par "QUESTION :" et n'ajoutez rien d'autre.} \\
            }
        }
    }
    & \\ \\

    \theutterance \stepcounter{utterance}  
    & \multicolumn{4}{p{0.6\linewidth}}{
        \cellcolor[rgb]{0.9,0.9,0.9}{
            \makecell[{{p{\linewidth}}}]{
                \texttt{\tiny{[P1$\rangle$GM]}}
                \texttt{QUESTION : Dans la troisième ligne de votre grille, les deux cases vides "▢" sont{-}elles suivies de trois "X" (en positions 3, 4 et 5) ?} \\
            }
        }
    }
    & & \\ \\

    \theutterance \stepcounter{utterance}  
    & & & \multicolumn{2}{p{0.3\linewidth}}{
        \cellcolor[rgb]{0.9,0.9,0.9}{
            \makecell[{{p{\linewidth}}}]{
                \texttt{\tiny{[GM$|$GM]}}
                \texttt{continue} \\
            }
        }
    }
    & & \\ \\

    \theutterance \stepcounter{utterance}  
    & & \multicolumn{4}{p{0.6\linewidth}}{
        \cellcolor[rgb]{0.9,0.9,0.9}{
            \makecell[{{p{\linewidth}}}]{
                \texttt{\tiny{[GM$\rangle$P2]}}
                \texttt{RÉPONSE : Oui, dans la deuxième ligne, les quatre cases vides "▢" sont suivies d'un "X" en cinquième position (bordure droite).} \\
\texttt{QUESTION : Dans la troisième ligne de votre grille, les deux cases vides "▢" sont{-}elles suivies de trois "X" (en positions 3, 4 et 5) ?Commencez votre réponse par "RÉPONSE :" et n'ajoutez rien d'autre.} \\
            }
        }
    }
    & \\ \\

    \theutterance \stepcounter{utterance}  
    & & & \multicolumn{4}{p{0.6\linewidth}}{
        \cellcolor[rgb]{0.9,0.9,0.9}{
            \makecell[{{p{\linewidth}}}]{
                \texttt{\tiny{[GM$\langle$P2]}}
                \texttt{RÉPONSE : Oui, dans la troisième ligne, les deux cases vides sont suivies de trois "X" en positions 3, 4 et 5.} \\
            }
        }
    }
     \\ \\

    \theutterance \stepcounter{utterance}  
    & & & \multicolumn{2}{p{0.3\linewidth}}{
        \cellcolor[rgb]{0.9,0.9,0.9}{
            \makecell[{{p{\linewidth}}}]{
                \texttt{\tiny{[GM$|$GM]}}
                \texttt{continue} \\
            }
        }
    }
    & & \\ \\

    \theutterance \stepcounter{utterance}  
    & & \multicolumn{4}{p{0.6\linewidth}}{
        \cellcolor[rgb]{0.9,0.9,0.9}{
            \makecell[{{p{\linewidth}}}]{
                \texttt{\tiny{[GM$\rangle$P2]}}
                \texttt{Posez maintenant une question afin de découvrir de nouveaux aspects de ma grille qui pourraient être différents de votre grille. Commencez par "QUESTION :" et n'ajoutez rien d'autre.} \\
            }
        }
    }
    & \\ \\

    \theutterance \stepcounter{utterance}  
    & & & \multicolumn{4}{p{0.6\linewidth}}{
        \cellcolor[rgb]{0.9,0.9,0.9}{
            \makecell[{{p{\linewidth}}}]{
                \texttt{\tiny{[GM$\langle$P2]}}
                \texttt{QUESTION : Les première, quatrième et cinquième lignes de votre grille sont{-}elles identiques à "X X X X X", "▢ ▢ ▢ ▢ X" et "▢ ▢ ▢ ▢ X" respectivement ?} \\
            }
        }
    }
     \\ \\

    \theutterance \stepcounter{utterance}  
    & & & \multicolumn{2}{p{0.3\linewidth}}{
        \cellcolor[rgb]{0.9,0.9,0.9}{
            \makecell[{{p{\linewidth}}}]{
                \texttt{\tiny{[GM$|$GM]}}
                \texttt{continue} \\
            }
        }
    }
    & & \\ \\

    \theutterance \stepcounter{utterance}  
    & & \multicolumn{4}{p{0.6\linewidth}}{
        \cellcolor[rgb]{0.9,0.9,0.9}{
            \makecell[{{p{\linewidth}}}]{
                \texttt{\tiny{[P1$\langle$GM]}}
                \texttt{RÉPONSE : Oui, dans la troisième ligne, les deux cases vides sont suivies de trois "X" en positions 3, 4 et 5.} \\
\texttt{QUESTION : Les première, quatrième et cinquième lignes de votre grille sont{-}elles identiques à "X X X X X", "▢ ▢ ▢ ▢ X" et "▢ ▢ ▢ ▢ X" respectivement ?Commencez votre réponse par "RÉPONSE :" et n'ajoutez rien d'autre.} \\
            }
        }
    }
    & \\ \\

    \theutterance \stepcounter{utterance}  
    & \multicolumn{4}{p{0.6\linewidth}}{
        \cellcolor[rgb]{0.9,0.9,0.9}{
            \makecell[{{p{\linewidth}}}]{
                \texttt{\tiny{[P1$\rangle$GM]}}
                \texttt{RÉPONSE : Oui, ces lignes sont identiques.} \\
            }
        }
    }
    & & \\ \\

    \theutterance \stepcounter{utterance}  
    & & & \multicolumn{2}{p{0.3\linewidth}}{
        \cellcolor[rgb]{0.9,0.9,0.9}{
            \makecell[{{p{\linewidth}}}]{
                \texttt{\tiny{[GM$|$GM]}}
                \texttt{continue} \\
            }
        }
    }
    & & \\ \\

    \theutterance \stepcounter{utterance}  
    & & \multicolumn{4}{p{0.6\linewidth}}{
        \cellcolor[rgb]{0.9,0.9,0.9}{
            \makecell[{{p{\linewidth}}}]{
                \texttt{\tiny{[P1$\langle$GM]}}
                \texttt{Posez maintenant une question afin de découvrir de nouveaux aspects de ma grille qui pourraient être différents de votre grille. Commencez par "QUESTION :" et n'ajoutez rien d'autre.} \\
            }
        }
    }
    & \\ \\

    \theutterance \stepcounter{utterance}  
    & \multicolumn{4}{p{0.6\linewidth}}{
        \cellcolor[rgb]{0.9,0.9,0.9}{
            \makecell[{{p{\linewidth}}}]{
                \texttt{\tiny{[P1$\rangle$GM]}}
                \texttt{QUESTION : Considérez{-}vous que nos grilles sont identiques, ou y a{-}t{-}il encore un détail que vous souhaitez vérifier ?} \\
            }
        }
    }
    & & \\ \\

    \theutterance \stepcounter{utterance}  
    & & & \multicolumn{2}{p{0.3\linewidth}}{
        \cellcolor[rgb]{0.9,0.9,0.9}{
            \makecell[{{p{\linewidth}}}]{
                \texttt{\tiny{[GM$|$GM]}}
                \texttt{continue} \\
            }
        }
    }
    & & \\ \\

    \theutterance \stepcounter{utterance}  
    & & \multicolumn{4}{p{0.6\linewidth}}{
        \cellcolor[rgb]{0.9,0.9,0.9}{
            \makecell[{{p{\linewidth}}}]{
                \texttt{\tiny{[GM$\rangle$P2]}}
                \texttt{RÉPONSE : Oui, ces lignes sont identiques.} \\
\texttt{QUESTION : Considérez{-}vous que nos grilles sont identiques, ou y a{-}t{-}il encore un détail que vous souhaitez vérifier ?Commencez votre réponse par "RÉPONSE :" et n'ajoutez rien d'autre.} \\
            }
        }
    }
    & \\ \\

    \theutterance \stepcounter{utterance}  
    & & & \multicolumn{4}{p{0.6\linewidth}}{
        \cellcolor[rgb]{0.9,0.9,0.9}{
            \makecell[{{p{\linewidth}}}]{
                \texttt{\tiny{[GM$\langle$P2]}}
                \texttt{RÉPONSE : Je considère que nos grilles sont identiques.} \\
            }
        }
    }
     \\ \\

    \theutterance \stepcounter{utterance}  
    & & & \multicolumn{2}{p{0.3\linewidth}}{
        \cellcolor[rgb]{0.9,0.9,0.9}{
            \makecell[{{p{\linewidth}}}]{
                \texttt{\tiny{[GM$|$GM]}}
                \texttt{continue} \\
            }
        }
    }
    & & \\ \\

    \theutterance \stepcounter{utterance}  
    & & \multicolumn{4}{p{0.6\linewidth}}{
        \cellcolor[rgb]{0.9,0.9,0.9}{
            \makecell[{{p{\linewidth}}}]{
                \texttt{\tiny{[GM$\rangle$P2]}}
                \texttt{Posez maintenant une question afin de découvrir de nouveaux aspects de ma grille qui pourraient être différents de votre grille. Commencez par "QUESTION :" et n'ajoutez rien d'autre.} \\
            }
        }
    }
    & \\ \\

    \theutterance \stepcounter{utterance}  
    & & & \multicolumn{4}{p{0.6\linewidth}}{
        \cellcolor[rgb]{0.9,0.9,0.9}{
            \makecell[{{p{\linewidth}}}]{
                \texttt{\tiny{[GM$\langle$P2]}}
                \texttt{QUESTION : Puisque nous avons confirmé que nos grilles sont identiques, n'y a{-}t{-}il plus aucun aspect à vérifier ou souhaitez{-}vous valider une dernière fois cette conclusion ?} \\
            }
        }
    }
     \\ \\

    \theutterance \stepcounter{utterance}  
    & & & \multicolumn{2}{p{0.3\linewidth}}{
        \cellcolor[rgb]{0.9,0.9,0.9}{
            \makecell[{{p{\linewidth}}}]{
                \texttt{\tiny{[GM$|$GM]}}
                \texttt{continue} \\
            }
        }
    }
    & & \\ \\

    \theutterance \stepcounter{utterance}  
    & & \multicolumn{4}{p{0.6\linewidth}}{
        \cellcolor[rgb]{0.9,0.9,0.9}{
            \makecell[{{p{\linewidth}}}]{
                \texttt{\tiny{[P1$\langle$GM]}}
                \texttt{RÉPONSE : Je considère que nos grilles sont identiques.} \\
\texttt{QUESTION : Puisque nous avons confirmé que nos grilles sont identiques, n'y a{-}t{-}il plus aucun aspect à vérifier ou souhaitez{-}vous valider une dernière fois cette conclusion ?Commencez votre réponse par "RÉPONSE :" et n'ajoutez rien d'autre.} \\
            }
        }
    }
    & \\ \\

    \theutterance \stepcounter{utterance}  
    & \multicolumn{4}{p{0.6\linewidth}}{
        \cellcolor[rgb]{0.9,0.9,0.9}{
            \makecell[{{p{\linewidth}}}]{
                \texttt{\tiny{[P1$\rangle$GM]}}
                \texttt{RÉPONSE : Je valide cette dernière fois notre conclusion : nos grilles sont identiques et il n'y a plus aucun aspect à vérifier.} \\
            }
        }
    }
    & & \\ \\

    \theutterance \stepcounter{utterance}  
    & & & \multicolumn{2}{p{0.3\linewidth}}{
        \cellcolor[rgb]{0.9,0.9,0.9}{
            \makecell[{{p{\linewidth}}}]{
                \texttt{\tiny{[GM$|$GM]}}
                \texttt{continue} \\
            }
        }
    }
    & & \\ \\

    \theutterance \stepcounter{utterance}  
    & & \multicolumn{4}{p{0.6\linewidth}}{
        \cellcolor[rgb]{0.9,0.9,0.9}{
            \makecell[{{p{\linewidth}}}]{
                \texttt{\tiny{[P1$\langle$GM]}}
                \texttt{Posez maintenant une question afin de découvrir de nouveaux aspects de ma grille qui pourraient être différents de votre grille. Commencez par "QUESTION :" et n'ajoutez rien d'autre.} \\
            }
        }
    }
    & \\ \\

    \theutterance \stepcounter{utterance}  
    & \multicolumn{4}{p{0.6\linewidth}}{
        \cellcolor[rgb]{0.9,0.9,0.9}{
            \makecell[{{p{\linewidth}}}]{
                \texttt{\tiny{[P1$\rangle$GM]}}
                \texttt{QUESTION : Puisque nous avons déjà confirmé l'identité de nos grilles sur tous les aspects, quel serait l'intérêt de chercher d'autres différences potentielles ?} \\
            }
        }
    }
    & & \\ \\

    \theutterance \stepcounter{utterance}  
    & & & \multicolumn{2}{p{0.3\linewidth}}{
        \cellcolor[rgb]{0.9,0.9,0.9}{
            \makecell[{{p{\linewidth}}}]{
                \texttt{\tiny{[GM$|$GM]}}
                \texttt{continue} \\
            }
        }
    }
    & & \\ \\

    \theutterance \stepcounter{utterance}  
    & & \multicolumn{4}{p{0.6\linewidth}}{
        \cellcolor[rgb]{0.9,0.9,0.9}{
            \makecell[{{p{\linewidth}}}]{
                \texttt{\tiny{[GM$\rangle$P2]}}
                \texttt{RÉPONSE : Je valide cette dernière fois notre conclusion : nos grilles sont identiques et il n'y a plus aucun aspect à vérifier.} \\
\texttt{QUESTION : Puisque nous avons déjà confirmé l'identité de nos grilles sur tous les aspects, quel serait l'intérêt de chercher d'autres différences potentielles ?Commencez votre réponse par "RÉPONSE :" et n'ajoutez rien d'autre.} \\
            }
        }
    }
    & \\ \\

    \theutterance \stepcounter{utterance}  
    & & & \multicolumn{4}{p{0.6\linewidth}}{
        \cellcolor[rgb]{0.9,0.9,0.9}{
            \makecell[{{p{\linewidth}}}]{
                \texttt{\tiny{[GM$\langle$P2]}}
                \texttt{RÉPONSE : Il n'y a aucun intérêt à chercher d'autres différences potentielles, nos grilles sont parfaitement identiques.} \\
            }
        }
    }
     \\ \\

    \theutterance \stepcounter{utterance}  
    & & & \multicolumn{2}{p{0.3\linewidth}}{
        \cellcolor[rgb]{0.9,0.9,0.9}{
            \makecell[{{p{\linewidth}}}]{
                \texttt{\tiny{[GM$|$GM]}}
                \texttt{continue} \\
            }
        }
    }
    & & \\ \\

    \theutterance \stepcounter{utterance}  
    & & \multicolumn{4}{p{0.6\linewidth}}{
        \cellcolor[rgb]{0.9,0.9,0.9}{
            \makecell[{{p{\linewidth}}}]{
                \texttt{\tiny{[GM$\rangle$P2]}}
                \texttt{Prenez maintenant une décision. Qu’en pensez{-}vous : votre grille et l’autre grille décrite sont{-}elles identiques ? Écrivez "DÉCISION : même grille" si vous pensez qu’il s’agit de la même grille ou "DÉCISION : grilles différentes" si vous pensez qu’il s’agit de grilles différentes. N’ajoutez rien d’autre.} \\
            }
        }
    }
    & \\ \\

    \theutterance \stepcounter{utterance}  
    & & & \multicolumn{4}{p{0.6\linewidth}}{
        \cellcolor[rgb]{0.9,0.9,0.9}{
            \makecell[{{p{\linewidth}}}]{
                \texttt{\tiny{[GM$\langle$P2]}}
                \texttt{DÉCISION : même grille} \\
            }
        }
    }
     \\ \\

    \theutterance \stepcounter{utterance}  
    & & & \multicolumn{2}{p{0.3\linewidth}}{
        \cellcolor[rgb]{0.9,0.9,0.9}{
            \makecell[{{p{\linewidth}}}]{
                \texttt{\tiny{[GM$|$GM]}}
                \texttt{continue} \\
            }
        }
    }
    & & \\ \\

    \theutterance \stepcounter{utterance}  
    & & & \multicolumn{2}{p{0.3\linewidth}}{
        \cellcolor[rgb]{0.9,0.9,0.9}{
            \makecell[{{p{\linewidth}}}]{
                \texttt{\tiny{[GM$|$GM]}}
                \texttt{success} \\
            }
        }
    }
    & & \\ \\

    \theutterance \stepcounter{utterance}  
    & & \multicolumn{4}{p{0.6\linewidth}}{
        \cellcolor[rgb]{0.9,0.9,0.9}{
            \makecell[{{p{\linewidth}}}]{
                \texttt{\tiny{[P1$\langle$GM]}}
                \texttt{RÉPONSE : Il n'y a aucun intérêt à chercher d'autres différences potentielles, nos grilles sont parfaitement identiques.} \\
\texttt{Prenez maintenant une décision. Qu’en pensez{-}vous : votre grille et l’autre grille décrite sont{-}elles identiques ? Écrivez "DÉCISION : même grille" si vous pensez qu’il s’agit de la même grille ou "DÉCISION : grilles différentes" si vous pensez qu’il s’agit de grilles différentes. N’ajoutez rien d’autre.} \\
            }
        }
    }
    & \\ \\

    \theutterance \stepcounter{utterance}  
    & \multicolumn{4}{p{0.6\linewidth}}{
        \cellcolor[rgb]{0.9,0.9,0.9}{
            \makecell[{{p{\linewidth}}}]{
                \texttt{\tiny{[P1$\rangle$GM]}}
                \texttt{DÉCISION : même grille} \\
            }
        }
    }
    & & \\ \\

    \theutterance \stepcounter{utterance}  
    & & & \multicolumn{2}{p{0.3\linewidth}}{
        \cellcolor[rgb]{0.9,0.9,0.9}{
            \makecell[{{p{\linewidth}}}]{
                \texttt{\tiny{[GM$|$GM]}}
                \texttt{continue} \\
            }
        }
    }
    & & \\ \\

    \theutterance \stepcounter{utterance}  
    & & & \multicolumn{2}{p{0.3\linewidth}}{
        \cellcolor[rgb]{0.9,0.9,0.9}{
            \makecell[{{p{\linewidth}}}]{
                \texttt{\tiny{[GM$|$GM]}}
                \texttt{success} \\
            }
        }
    }
    & & \\ \\

\end{supertabular}
}

\end{document}
